\section*{Chapitre \ref{ch:s2:commodity-spreads} --- Spreads sur matières premières}

\begin{solution}{exo:s2:commodity-spreads:1}
$(2 \times 2.50 + 3.00) \times 42 / 3 - 80 = 112 - 80 = \$32$ par baril.
\end{solution}

\begin{solution}{exo:s2:commodity-spreads:2}
$50 - 7 \times 3 = \$29$ par mégawattheure~; à un taux de chaleur de 7.5, $50 - 22.5 = \$27.50$. Une centrale moins efficace gagne moins aux mêmes prix.
\end{solution}

\begin{solution}{exo:s2:commodity-spreads:3}
$0.022 \times 300 + 0.11 \times 50 - 10 = 6.60 + 5.50 - 10 = \$2.10$ par boisseau (le pétrole à 50 cents la livre vaut 0.50~\$ la livre, et $0.11$ vaut 11
livres en dollars par cent).
\end{solution}

\begin{solution}{exo:s2:commodity-spreads:4}
Les prix des produits suivent le brut avec un décalage~: les marchés de détail et de gros des produits s'ajustent en quelques jours ou semaines, si bien qu'un bond du brut resserre d'abord le crack spread et
une baisse l'élargit d'abord. Les variations quotidiennes du crack spread 3-2-1 réel ont une corrélation de $-0.38$ avec celles du WTI~; celle du crack spread synthétique, $-0.36$, a été injectée.
\end{solution}

\begin{solution}{exo:s2:commodity-spreads:5}
Le Brent est évalué à la clôture de Londres et le WTI à celle de New York, si bien que le spread d'une journée combine des prix pris à des heures différentes~; un mouvement du pétrole survenant après la
clôture de Londres apparaît dans le WTI le jour même et dans le Brent le lendemain, et la variation du spread s'inverse. Les variations quotidiennes ont une autocorrélation de $-0.36$. Une règle
qui s'exécute à la clôture de la journée capte cette inversion, que personne ne pourrait négocier à ces prix~; passée avec un jour de retard, son ratio de Sharpe tombe de 0.85
à 0.32.
\end{solution}

\begin{solution}{exo:s2:commodity-spreads:6}
Capacité de trituration~: quand le crush spread est large, les transformateurs achètent des graines de soja, vendent du tourteau et de l'huile, et triturent davantage~; quand il est étroit, ils ralentissent. Leur
arbitrage ramène le spread vers le coût de transformation, comme Simon l'a constaté pour 1985--1995.
\end{solution}

\begin{solution}{exo:s2:commodity-spreads:7}
Ratios de Sharpe de 0.85, 1.05 et 0.65, avec une présence sur le marché de 73~\%, 56~\% et 31~\% des jours. Choisir 1.5 après avoir vu les trois est un petit test multiple~:
le 1.05 publié est le meilleur de trois essais et surestime ce que la règle gagnerait ensuite (le ratio de Sharpe déflaté du livre~4 s'applique). La bande doit
être fixée avant le test, ou choisie sur des données antérieures à la période d'évaluation.
\end{solution}

\begin{solution}{exo:s2:commodity-spreads:8}
Le spread ne revient au coût de transport que tant que la capacité de transport est disponible. Sur le marché synthétique, un goulet d'étranglement sur un pipeline a élargi le spread de
15~\$ en six mois, et le fade a perdu 14.1~\$ pendant ces mois après avoir gagné 35.3~\$ en quatre ans~; doubler la position double cette perte.
L'EIA a attribué le record du spread Brent-WTI de 2011 à des goulets d'étranglement près de Cushing. Jouer le fade sur un \omterm{def:s2:commodity-spreads:location}{spread géographique} exige une vue sur la capacité, et un stop
pour le cas où elle est saturée.
\end{solution}

\begin{solution}{pb:s2:commodity-spreads:1}
\begin{enumerate}
\item Un \omterm{def:s2:commodity-spreads:processing}{trade de marge de transformation} détient un intrant face à ses extrants dans le ratio du procédé, pour en tirer la marge~; un \omterm{def:s2:commodity-spreads:location}{spread géographique} est l'écart de prix d'une même
matière première entre deux lieux~; un \omterm{def:s2:commodity-spreads:quality}{spread de qualité} est la différence entre deux qualités.
\item Capacité de transformation et son coût (raffinage, production d'électricité, trituration)~; capacité et coût du transport~; coût de la valorisation de la qualité inférieure.
\item Les contrats à terme sur brut, essence et fioul domestique sont cointégrés avec des spreads stationnaires, et l'arbitrage de risque sur trois spreads a été historiquement
rentable, bien que peut-être plus aujourd'hui.
\item Les écarts par rapport à un équilibre saisonnier et tendanciel étaient transitoires, avec un retour vers la moyenne récente sur cinq jours, et des règles simulées auraient
été rentables.
\item Moyenne 21.30~\$, demi-vie d'environ 51 jours de bourse, corrélation des variations quotidiennes avec celles du WTI de $-0.38$.
\item En 2011, le Brent a valu en moyenne environ 111~\$ et le WTI environ 95~\$, un spread record que l'EIA a attribué à des goulets d'étranglement près de Cushing~; le 20 avril 2020, le prix spot du WTI
était de $-\$36.98$ et le Brent se situait 54.34~\$ au-dessus de celui-ci.
\item Les deux prix sont évalués à des clôtures différentes~; les variations ont une autocorrélation de $-0.36$, et une règle à même clôture capte une inversion que
personne ne pourrait négocier~: le ratio de Sharpe du fade tombe de 0.85 à 0.32 avec un jour de retard.
\item Il a progressé 14 années sur 20, de 6.17~\$ en moyenne, avec une statistique $t$ de 2.8, et la pire année à $-\$11.97$.
\item Une moyenne de 20~\$, un cycle saisonnier de 5~\$ culminant fin avril, un écart qui revient vers sa moyenne avec une demi-vie de 40 jours, et un décalage injecté par rapport au brut.
\item Fade sur le crack spread, 15.4~\$ par an, ratio de Sharpe 1.05, corrélation $-0.12$~; saison retirée 6.2, 0.44, $-0.07$~; février--avril 7.4, 0.78, $-0.15$~;
\omterm{def:s2:commodity-spreads:location}{spread géographique} 3.8, 0.48, 0.01~; \omterm{def:s2:commodity-spreads:quality}{spread de qualité} 2.6, 0.98, 0.01~; combinaison 1.69, $-0.13$.
\item La règle simple parie aussi bien contre la saison que contre l'écart~; sans la saison, il ne reste que l'écart.
\item Il a gagné 35.3~\$ sur quatre ans, perdu 14.1~\$ dans les six premiers mois du goulet d'étranglement et 9.3~\$ sur toute sa durée de vie, puis gagné 8.4~\$ après.
\item Le fade synthétique sur le crack spread affiche un ratio de Sharpe de 1.05, avec une corrélation de $-0.12$ aux variations quotidiennes du brut~; sur les données spot de l'EIA, 0.48 avec
$-0.02$ (0.35 sans avril 2020, 0.24 avec un jour de retard).
\item Deux de ses quatre règles sont un seul pari sur la saison du crack spread~; les effets injectés sont forts~; et le risque de la combinaison est mesuré sur un échantillon
comportant un seul goulet d'étranglement.
\item Un retour à la moyenne qui dépend d'un lien physique pouvant se rompre~; des jambes ayant chacune leur propre marge, liquidité et échéance~; des ratios qui sont des conventions.
\item Sur contrats à terme et non sur le spot, à clôtures synchrones, avec des paramètres fixés avant le test et la saison estimée sur les seules années passées.
\item Quand la capacité de transport est saturée~: le spread traduit alors la rareté du pipeline, et non un écart.
\item Le spark spread avec une vue sur le taux de chaleur~: il exige un modèle du parc de centrales.
\item Les deux approches opèrent sur de nombreux marchés à partir de signaux simples et à risque égal~; les spreads éliminent le niveau de chaque matière première, que le macro systématique conserve.
\item Sur l'activité physique qui relie les jambes~: la marge du raffinage, de la production d'électricité, de la trituration ou du transport de la matière première.
\end{enumerate}
\end{solution}

\begin{solution}{iq:s2:commodity-spreads:1}
Trois barils de brut contre deux d'essence et un de distillat~: une marge de raffinage grossière. Les raffineurs la couvrent en vendant les produits et en achetant du brut
à terme~; les compagnies aériennes et les acheteurs de produits couvrent l'autre côté.
\end{solution}

\begin{solution}{iq:s2:commodity-spreads:2}
Vérifier que les jambes sont enregistrées à la même heure~; tester la règle passée avec un jour de retard~; utiliser les contrats à terme négociés et non les évaluations~; examiner
l'autocorrélation des variations quotidiennes (une forte autocorrélation négative signale des prix asynchrones ou bruités)~; et vérifier que le retour à la moyenne survit aux coûts
et aux périodes hors échantillon.
\end{solution}

\begin{solution}{iq:s2:commodity-spreads:3}
Seulement après s'être demandé pourquoi~: si les pipelines ou la capacité d'exportation sont saturés, le spread traduit un transport rare et peut continuer à s'élargir, comme en 2011~; si
il s'agit d'une panne temporaire ou d'un serrage de livraison, un fade assorti d'un stop peut rapporter. Selon l'EIA, les variations à court terme viennent de la capacité des pipelines ou
de la production.
\end{solution}

\begin{solution}{iq:s2:commodity-spreads:4}
Des limites sur le risque de chaque spread et sur la position directionnelle et la liquidité de chaque jambe~; des tests de résistance qui brisent le lien de chaque spread (goulets d'étranglement, pannes,
serrages de livraison) et un scénario commun à tous~; des limites à l'approche des échéances et de la livraison.
\end{solution}

\begin{solution}{iq:s2:commodity-spreads:5}
Conserver chaque jambe comme une série de contrats à terme, avec son propre calendrier de roll~; construire le spread à partir des jambes à des horodatages synchrones~; rouler les jambes
ensemble, ou enregistrer le décalage de roll~; calculer les coûts par jambe, et la marge par jambe ou par spread reconnu.
\end{solution}

\begin{solution}{iq:s2:commodity-spreads:6}
Découper le jour $d$ en la partie antérieure à l'heure d'enregistrement la plus tôt et la dernière fraction $f$. Le spread est l'accroissement du prix sur le dernier $f$ de la
journée, $S_d = W_d$. Ainsi $\Delta S_d = W_d - W_{d-1}$ et $\Delta S_{d+1} = W_{d+1} - W_d$, avec des accroissements indépendants de variance $f\sigma^2$~: la
covariance vaut $-f\sigma^2$ et la variance $2f\sigma^2$, si bien que l'autocorrélation vaut $-1/2$ quelle que soit $f$. Si le spread vrai bouge aussi, ses propres variations
ajoutent de la variance sans cette covariance, et l'autocorrélation se situe entre $-1/2$ et zéro, comme la valeur réelle $-0.36$.
\end{solution}
